\chapter{One Circuit, Four Views}
\section{The Recurring Framework}
The supplied draft's core idea is the organizing device for the rest of this
book. A circuit can be viewed as:
\begin{enumerate}
\item a physical network governed by conservation laws;
\item a differential equation;
\item a state-space model;
\item a phasor or frequency-response model;
\item an energy-storage and dissipation process.
\end{enumerate}
These views should agree. If they do not, at least one sign convention,
reference direction, or assumption needs attention.

\section{Series RC Circuit}
Take a source \(u(t)\) driving a resistor \(R\) in series with a capacitor \(C\).
Let \(v_C(t)\) be the capacitor voltage and \(i(t)\) the series current.

\begin{figure}[htbp]
\centering
\begin{circuitikz}
\draw (0,0) to[sV,l=\(u(t)\)] (0,2)
      to[R,l=\(R\),i=\(i(t)\)] (3,2)
      to[C,l=\(C\),v=\(v_C(t)\)] (3,0)
      -- (0,0);
\end{circuitikz}
\caption{Series RC circuit with output across the capacitor.}
\end{figure}

\subsection{Physical and Energy View}
The capacitor stores electric-field energy
\[
E_C=\frac{1}{2}Cv_C^2.
\]
The resistor dissipates energy, so with no input the system relaxes toward
equilibrium.

\subsection{Differential Equation}
KVL gives
\[
u(t)=v_R(t)+v_C(t)=Ri(t)+v_C(t),
\qquad
i(t)=C\frac{\dd v_C}{\dd t}.
\]
Therefore
\[
RC\frac{\dd v_C}{\dd t}+v_C=u(t).
\]

\subsection{State-Space Form}
With \(x=v_C\),
\[
\dot{x}=-\frac{1}{RC}x+\frac{1}{RC}u.
\]
There is only one state because there is only one independent storage element.
The stable eigenvalue is \(-1/(RC)\).

\subsection{Frequency-Response Form}
With zero initial conditions and sinusoidal steady state,
\[
G(s)=\frac{V_C(s)}{U(s)}=\frac{1}{1+sRC},
\qquad
G(\jj\omega)=\frac{1}{1+\jj\omega RC}.
\]
High frequency makes the capacitor impedance small, so the output across the
capacitor is attenuated.

\section{Series RL Circuit}
Now replace the capacitor with an inductor \(L\). KVL gives
\[
u(t)=Ri(t)+L\frac{\dd i}{\dd t},
\]
or
\[
L\frac{\dd i}{\dd t}+Ri=u(t).
\]
With \(x=i\),
\[
\dot{x}=-\frac{R}{L}x+\frac{1}{L}u,
\qquad
\tau=\frac{L}{R}.
\]
The inductor stores magnetic-field energy \(E_L=\frac{1}{2}Li^2\).

\section{Series RLC Circuit}
For a source driving \(R\), \(L\), and \(C\) in series,
\[
LC\frac{\dd^2 v_C}{\dd t^2}
+RC\frac{\dd v_C}{\dd t}+v_C=u(t).
\]
Using \(x_1=v_C\) and \(x_2=i\),
\[
\dot{x}_1=\frac{1}{C}x_2,
\qquad
\dot{x}_2=-\frac{1}{L}x_1-\frac{R}{L}x_2+\frac{1}{L}u.
\]
The phasor impedance is
\[
Z=R+\jj\left(\omega L-\frac{1}{\omega C}\right).
\]
Resonance occurs when \(\omega L=1/(\omega C)\).

\begin{energybox}
For the series RLC circuit,
\[
\mathcal{H}=\frac{1}{2}Cv_C^2+\frac{1}{2}Li^2.
\]
With \(u=v_R+v_L+v_C\), \(i=C\dot v_C\), and \(v_L=L\dot i\),
\[
\frac{\dd\mathcal{H}}{\dd t}
=Cv_C\dot v_C+Li\dot i
=v_Ci+iv_L
=i(u-v_R)
=ui-Ri^2.
\]
Stored energy changes at input power minus resistive dissipation.
\end{energybox}

\begin{controlsbox}
The RLC circuit is the electrical analog of a mass-spring-damper system.
Eigenvalues are natural modes. Higher resistance gives stronger damping; lower
loss gives more persistent oscillation.
\end{controlsbox}
